% GENERATED FILE. Edit canonical lessons, then run npm run generate:books.
% canonical-source-hash: fnv1a64:dfed931a826cda0b
% canonical-lessons: BN-C387-poth, BN-C387-gontobyo, BN-C387-sophor, BN-C387-jatayat-kora, BN-C387-phirti

\chapter{Ways and Destinations}
\label{ch:bn-ways-and-destinations}
\hlchaptermodality{387}

\begin{chapteropening}
\textbf{By the end of this chapter:} \emph{I can talk about the way, where a trip ends and the journey back.}

Complete the last lesson of chapter 387: I can talk about the way, where a trip ends and the journey back.
\end{chapteropening}

\section[\texorpdfstring{\bn{পথ} (pôth)}{পথ (pôth)}]{\bn{পথ} (pôth) — a way; a path}

\label{lesson:BN-C387-poth}



\begin{quote}
Before the new one: say the Bengali for to unload, then the Bengali for to carry (a load).
\end{quote}

\subsection*{You'll want to know: \bn{পথ}}
\textbf{\bn{পথ}} — \emph{pôth} — \textquotedblleft{}a way; a path\textquotedblright{}.

\begin{quote}
\bn{গ্রামের} \bn{পথে} \bn{আমরা} \bn{একটা} \bn{পুকুর} \bn{দেখলাম।}
\end{quote}

\begin{quote}
\emph{grāmer pôthe āmrā ekṭā pukur dekhlām.}
\end{quote}

\begin{quote}
On the village path we saw a pond.
\end{quote}

\bn{পথে}, on the way: \bn{পথে} \bn{দেখা} \bn{হল}, we met on the way. \bn{রাস্তা} is the road itself; \bn{পথ} is also the way you take.

\begin{cousinweb}
\textbf{\bn{পথ}} (\emph{pôth}) is Sanskrit \emph{patha}, \textquotedblleft{}a way, a path, a road\textquotedblright{} (Monier-Williams), the stem Sanskrit uses at the end of compounds, taken over whole.
\end{cousinweb}

\begin{grammarlens}[title={What you've built}]
One more word for getting around.
\end{grammarlens}

\subsection*{Your turn}
\begin{itemize}
\raggedright
  \item \emph{Say it:} \emph{pôth}
  \item \emph{Say it:} \emph{pôth}, once more
  \item \emph{Say it:} \emph{grāmer pôthe āmrā ekṭā pukur dekhlām.} — On the village path we saw a pond.
  \item \emph{Recall:} say the Bengali for to unload, then the Bengali for to carry (a load), then say \emph{pôth} again
\end{itemize}

\begin{tcolorbox}[breakable,colback=teal!4,colframe=teal!35!black,title={Before you move on}]
What does \bn{পথ} mean? (\textquotedblleft{}A way; a path\textquotedblright{}.) Say it once more.
\end{tcolorbox}
\section[\texorpdfstring{\bn{গন্তব্য} (gôntôbyô)}{গন্তব্য (gôntôbyô)}]{\bn{গন্তব্য} (gôntôbyô) — a destination}

\label{lesson:BN-C387-gontobyo}



\begin{quote}
Before the new one: say the Bengali for to carry (a load), then the Bengali for a path.
\end{quote}

\subsection*{You'll want to know: \bn{গন্তব্য}}
\textbf{\bn{গন্তব্য}} — \emph{gôntôbyô} — \textquotedblleft{}a destination\textquotedblright{}.

\begin{quote}
\bn{আমাদের} \bn{গন্তব্য} \bn{একটা} \bn{ছোট} \bn{গ্রাম।}
\end{quote}

\begin{quote}
\emph{āmāder gôntôbyô ekṭā chhoṭo grām.}
\end{quote}

\begin{quote}
Our destination is a small village.
\end{quote}

A word of signs and tickets as much as of talk: \bn{গন্তব্য} \bn{স্টেশন}, the destination station.

\begin{cousinweb}
\textbf{\bn{গন্তব্য}} (\emph{gôntôbyô}) is Sanskrit \emph{gantavya}, \textquotedblleft{}to be gone to, to be reached\textquotedblright{} (Monier-Williams), a form of \emph{gam}, to go: the place you must go to.
\end{cousinweb}

\begin{grammarlens}[title={What you've built}]
One more word for getting around.
\end{grammarlens}

\subsection*{Your turn}
\begin{itemize}
\raggedright
  \item \emph{Say it:} \emph{gôntôbyô}
  \item \emph{Say it:} \emph{gôntôbyô}, once more
  \item \emph{Say it:} \emph{āmāder gôntôbyô ekṭā chhoṭo grām.} — Our destination is a small village.
  \item \emph{Recall:} say the Bengali for to carry (a load), then the Bengali for a path, then say \emph{gôntôbyô} again
\end{itemize}

\begin{tcolorbox}[breakable,colback=teal!4,colframe=teal!35!black,title={Before you move on}]
What does \bn{গন্তব্য} mean? (\textquotedblleft{}A destination\textquotedblright{}.) Say it once more.
\end{tcolorbox}
\section[\texorpdfstring{\bn{সফর} (sôphôr)}{সফর (sôphôr)}]{\bn{সফর} (sôphôr) — a trip; a tour}

\label{lesson:BN-C387-sophor}



\begin{quote}
Before the new one: say the Bengali for a path, then the Bengali for a destination.
\end{quote}

\subsection*{You'll want to know: \bn{সফর}}
\textbf{\bn{সফর}} — \emph{sôphôr} — \textquotedblleft{}a trip; a tour\textquotedblright{}.

\begin{quote}
\bn{আমাদের} \bn{সফর} \bn{সাত} \bn{দিনের} \bn{ছিল।}
\end{quote}

\begin{quote}
\emph{āmāder sôphôr sāt diner chhilô.}
\end{quote}

\begin{quote}
Our trip lasted seven days.
\end{quote}

Often a trip with a purpose: an official tour, a team's tour abroad.

\begin{cousinweb}
\textbf{\bn{সফর}} (\emph{sôphôr}) is usually traced to Arabic \emph{safar}, \textquotedblleft{}a journey\textquotedblright{}, which reached Bengali through Persian.
\end{cousinweb}

\begin{grammarlens}[title={What you've built}]
One more word for getting around.
\end{grammarlens}

\subsection*{Your turn}
\begin{itemize}
\raggedright
  \item \emph{Say it:} \emph{sôphôr}
  \item \emph{Say it:} \emph{sôphôr}, once more
  \item \emph{Say it:} \emph{āmāder sôphôr sāt diner chhilô.} — Our trip lasted seven days.
  \item \emph{Recall:} say the Bengali for a path, then the Bengali for a destination, then say \emph{sôphôr} again
\end{itemize}

\begin{tcolorbox}[breakable,colback=teal!4,colframe=teal!35!black,title={Before you move on}]
What does \bn{সফর} mean? (\textquotedblleft{}A trip; a tour\textquotedblright{}.) Say it once more.
\end{tcolorbox}
\section[\texorpdfstring{\bn{যাতায়াত} \bn{করা} (jātāyāt kôrā)}{যাতায়াত করা (jātāyāt kôrā)}]{\bn{যাতায়াত} \bn{করা} (jātāyāt kôrā) — to commute; to go back and forth}

\label{lesson:BN-C387-jatayat-kora}



\begin{quote}
Before the new one: say the Bengali for a destination, then the Bengali for a trip.
\end{quote}

\subsection*{You'll want to know: \bn{যাতায়াত} \bn{করা}}
\textbf{\bn{যাতায়াত} \bn{করা}} — \emph{jātāyāt kôrā} — \textquotedblleft{}to commute; to go back and forth\textquotedblright{}.

\begin{quote}
\bn{মীরা} \bn{রোজ} \bn{বাসে} \bn{যাতায়াত} \bn{করে।}
\end{quote}

\begin{quote}
\emph{mīrā roj bāse jātāyāt kôre.}
\end{quote}

\begin{quote}
Mira commutes by bus every day.
\end{quote}

Past (he/she …): \bn{তিনি} \bn{ট্রেনে} \bn{যাতায়াত} \bn{করতেন}, \emph{tini ṭrene jātāyāt kôrten}.

The noun alone is \bn{যাতায়াত}, the coming and going: \bn{এই} \bn{রাস্তায়} \bn{যাতায়াত}, the traffic on this road.

\begin{cousinweb}
\textbf{\bn{যাতায়াত} \bn{করা}} (\emph{jātāyāt kôrā}) joins Sanskrit \emph{yātāyāta}, \textquotedblleft{}going and coming\textquotedblright{} (Monier-Williams), a pairing of \emph{yāta}, gone, and \emph{āyāta}, come, to \bn{করা}, to do.
\end{cousinweb}

\begin{grammarlens}[title={What you've built}]
One more word for getting around.
\end{grammarlens}

\subsection*{Your turn}
\begin{itemize}
\raggedright
  \item \emph{Say it:} \emph{jātāyāt kôrā}
  \item \emph{Say it:} \emph{jātāyāt kôrā}, once more
  \item \emph{Say it:} \emph{mīrā roj bāse jātāyāt kôre.} — Mira commutes by bus every day.
  \item \emph{Recall:} say the Bengali for a destination, then the Bengali for a trip, then say \emph{jātāyāt kôrā} again
\end{itemize}

\begin{tcolorbox}[breakable,colback=teal!4,colframe=teal!35!black,title={Before you move on}]
What does \bn{যাতায়াত} \bn{করা} mean? (\textquotedblleft{}To commute; to go back and forth\textquotedblright{}.) Say it once more.
\end{tcolorbox}
\section[\texorpdfstring{\bn{ফিরতি} (phirti)}{ফিরতি (phirti)}]{\bn{ফিরতি} (phirti) — return (of a journey); on the way back}

\label{lesson:BN-C387-phirti}



\begin{quote}
Before the new one: say the Bengali for a trip, then the Bengali for to commute.
\end{quote}

\subsection*{You'll want to know: \bn{ফিরতি}}
\textbf{\bn{ফিরতি}} — \emph{phirti} — \textquotedblleft{}return (of a journey); on the way back\textquotedblright{}.

\begin{quote}
\bn{ফিরতি} \bn{পথে} \bn{আমরা} \bn{মামার} \bn{বাড়িতে} \bn{থামলাম।}
\end{quote}

\begin{quote}
\emph{phirti pôthe āmrā māmār bāṛite thāmlām.}
\end{quote}

\begin{quote}
On the way back we stopped at my uncle's house.
\end{quote}

It stands before its noun: \bn{ফিরতি} \bn{টিকিট}, a return ticket; \bn{ফিরতি} \bn{ট্রেন}, the train back.

\begin{cousinweb}
\textbf{\bn{ফিরতি}} (\emph{phirti}) is built on the stem \bn{ফির}- of \textbf{\bn{ফেরা}}, \textquotedblleft{}to return\textquotedblright{}, which CDIAL lists as Bengali \emph{phirā}, \textquotedblleft{}to turn\textquotedblright{}, under the reconstructed \emph{phirati}, \textquotedblleft{}moves, wanders, turns\textquotedblright{} (CDIAL 9078).
\end{cousinweb}

\begin{grammarlens}[title={What you've built}]
One more word for getting around.
\end{grammarlens}

\subsection*{Your turn}
\begin{itemize}
\raggedright
  \item \emph{Say it:} \emph{phirti}
  \item \emph{Say it:} \emph{phirti}, once more
  \item \emph{Say it:} \emph{phirti pôthe āmrā māmār bāṛite thāmlām.} — On the way back we stopped at my uncle's house.
  \item \emph{Recall:} say the Bengali for a trip, then the Bengali for to commute, then say \emph{phirti} again
\end{itemize}

\begin{tcolorbox}[breakable,colback=teal!4,colframe=teal!35!black,title={Before you move on}]
What does \bn{ফিরতি} mean? (\textquotedblleft{}Return (of a journey); on the way back\textquotedblright{}.) Say it once more.
\end{tcolorbox}
